\uuid{lhQg}
\titre{Trigonométrie}
\auteur{bodin ; hanani ; mzari}
\organisation{exo7}
\licence{CC BY-NC-SA 4.0}
\datecreate{2020-01-20}
\niveau{L1}
\module{Algèbre}
\chapitre{Nombres complexes}
\sousChapitre{Forme cartésienne, forme polaire}
\difficulte{1}
\theme{Écritures algébrique et géométrique, trigonométrie}
\exo7qcmid{lille-1.nombres-complexes.005}

\contenu{
\begin{question}
\qclassification{104.01}
\qtags{motcle=trigonométrie}
Soit $\theta$ un réel.  Quelles sont les assertions vraies ?
\begin{answers}
    \good{$\cos^2\theta= \frac{1+\cos(2\theta)}{2}$}
    \bad{ $\cos^2\theta= \frac{1-\cos(2\theta)}{2}$}

    \good{$\sin^2\theta= \frac{1-\cos(2\theta)}{2}$}

    \bad{$\sin^2\theta= \frac{1+\cos(2\theta)}{2}$}
\end{answers}
\begin{explanations}
On peut appliquer les formules d'Euler, ou utiliser la formule d'addition du cosinus.

\end{explanations}
\end{question}
}
